\documentclass[11pt,a4paper]{article}
% =====================================================================
%  Shared preamble for the Part V lecture notes (lectures 19-22).
%
%  These four lectures sit outside Geron, so the notes ARE the primary
%  source and are examined on the same terms as the textbook parts.
%  Everything here is chosen to make that readable on paper: one column,
%  generous measure, and the same six-colour palette as the slides so a
%  student moving between the two is not re-learning what red means.
%
%  Build with pdflatex (twice, for the toc and the refs):
%      cd notes && latexmk -pdf lecture-19.tex
%  or  make -C notes
% =====================================================================
\usepackage[T1]{fontenc}
\usepackage[utf8]{inputenc}
\usepackage{newtxtext,newtxmath}      % Times-like: prints well, reads well
\usepackage[scaled=0.86]{inconsolata} % code, at a size that sits beside Times
\usepackage{microtype}
% newtx defines \Bbbk, \openbox, \proof and \endproof; amssymb and amsthm
% then redefine them and abort. Freeing the four names before those two
% packages load is the standard fix, and it costs nothing here: we use
% amsthm only for \newtheorem.
\let\Bbbk\relax
\usepackage{amsmath,amssymb}
\let\openbox\relax
\let\proof\relax
\let\endproof\relax
\usepackage{amsthm}
\usepackage{booktabs}
\usepackage{enumitem}
\usepackage{xcolor}
\usepackage{tcolorbox}
\usepackage{titlesec}
\usepackage{graphicx}
\usepackage[hidelinks]{hyperref}
\tcbuselibrary{skins,breakable}

% ---- the course palette, from assets/css/custom.css -------------------
\definecolor{aimlprimary}{HTML}{0B3D62}   % deep blue  - structure
\definecolor{aimlaccent} {HTML}{C0392B}   % brick red  - failures
\definecolor{aimlsuccess}{HTML}{14663A}   % green      - repairs
\definecolor{aimlmath}   {HTML}{6C3483}   % purple     - the derivation
\definecolor{aimlmuted}  {HTML}{4B5563}
\definecolor{aimlrule}   {HTML}{B0BCC7}
\definecolor{aimlsoft}   {HTML}{F4F7F9}

\usepackage[a4paper,margin=32mm,top=30mm,bottom=30mm]{geometry}
\setlength{\parskip}{0.45\baselineskip}
\setlength{\parindent}{0pt}
\linespread{1.06}

% ---- headings --------------------------------------------------------
\titleformat{\section}{\Large\bfseries\color{aimlprimary}}{\thesection}{0.7em}{}
\titleformat{\subsection}{\large\bfseries\color{aimlprimary}}{\thesubsection}{0.7em}{}
\titleformat{\subsubsection}{\bfseries\color{aimlmuted}}{}{0em}{}
\titlespacing*{\section}{0pt}{1.6\baselineskip}{0.5\baselineskip}
\titlespacing*{\subsection}{0pt}{1.2\baselineskip}{0.35\baselineskip}

% ---- boxes -----------------------------------------------------------
% One box per role, and the role is the colour. A student who has seen the
% slides already knows what each one means before reading a word of it.
\newtcolorbox{keybox}[1][]{
  breakable, enhanced, colback=aimlsoft, colframe=aimlprimary,
  boxrule=0.4pt, left=8pt, right=8pt, top=6pt, bottom=6pt,
  fonttitle=\bfseries\small, coltitle=white, colbacktitle=aimlprimary,
  attach boxed title to top left={yshift=-2mm, xshift=4mm},
  boxed title style={boxrule=0pt, arc=1pt}, #1}

\newtcolorbox{failbox}[1][]{
  breakable, enhanced, colback=white, colframe=aimlaccent,
  boxrule=0.9pt, left=8pt, right=8pt, top=6pt, bottom=6pt,
  fonttitle=\bfseries\small, coltitle=white, colbacktitle=aimlaccent,
  attach boxed title to top left={yshift=-2mm, xshift=4mm},
  boxed title style={boxrule=0pt, arc=1pt}, #1}

\newtcolorbox{derivbox}[1][]{
  breakable, enhanced, colback=white, colframe=aimlmath,
  boxrule=0.6pt, left=8pt, right=8pt, top=6pt, bottom=6pt,
  fonttitle=\bfseries\small, coltitle=white, colbacktitle=aimlmath,
  attach boxed title to top left={yshift=-2mm, xshift=4mm},
  boxed title style={boxrule=0pt, arc=1pt}, #1}

% An exercise. Deliberately the same shape as the notebook's `try` field:
% one change, and what should happen to the output.
\newtcolorbox{trybox}[1][]{
  breakable, enhanced, colback=white, colframe=aimlsuccess,
  boxrule=0.5pt, left=8pt, right=8pt, top=6pt, bottom=6pt,
  fonttitle=\bfseries\small, coltitle=white, colbacktitle=aimlsuccess,
  attach boxed title to top left={yshift=-2mm, xshift=4mm},
  boxed title style={boxrule=0pt, arc=1pt}, #1}

% ---- small semantics -------------------------------------------------
\newcommand{\scope}[1]{\textcolor{aimlmuted}{\small #1}}
\newcommand{\term}[1]{\textbf{#1}}
\newcommand{\code}[1]{\texttt{\small #1}}
\newcommand{\lecture}[1]{Lecture~#1}

\newtheorem{definition}{Definition}[section]
\newtheorem{proposition}[definition]{Proposition}

% ---- title block -----------------------------------------------------
% \notestitle{19}{Information retrieval: the lexical foundation}{SciFact (BEIR)}
\newcommand{\notestitle}[3]{%
  \thispagestyle{empty}
  \begin{flushleft}
    {\color{aimlmuted}\small\sffamily
     APPLICATIONS OF MACHINE LEARNING \quad$\cdot$\quad
     BSc Mathematics of Artificial Intelligence}\\[2mm]
    {\color{aimlrule}\rule{\textwidth}{0.8pt}}\\[4mm]
    {\color{aimlmuted}\large Lecture #1 \quad$\cdot$\quad Part V \quad$\cdot$\quad
     Extended notes}\\[3mm]
    {\Huge\bfseries\color{aimlprimary}\baselineskip=1.15\baselineskip #2\par}\vspace{4mm}
    {\color{aimlmuted}\large Dataset: #3}\\[3mm]
    {\color{aimlrule}\rule{\textwidth}{0.8pt}}
  \end{flushleft}
  \vspace{2mm}
  \begin{keybox}[title={Why these notes exist}]
    Lectures 19--22 sit outside G\'eron, the course's primary text. For those
    four lectures \emph{these notes are the primary source}, and they are
    examined on exactly the same terms as the textbook parts. Everything in
    the slide deck for this lecture is here, with the arguments written out at
    length rather than compressed onto a slide; the numbers are the ones the
    accompanying notebook prints, and every one of them is a measurement on
    the corpus named above rather than a figure quoted from a paper.
  \end{keybox}
  \vspace{1mm}
  {\small\tableofcontents}
  \vspace{2mm}
  \noindent{\color{aimlrule}\rule{\textwidth}{0.4pt}}
  \vspace{1mm}
}


\begin{document}
\notestitlefor{24}{VI}{Generation, retrieval-augmented systems, and where this leaves you}{COCO and the Part V corpora}{Chapters 15--16}

\section{Two weak spots, both about entries}

Lecture 23 built a catalogue you can search with a sentence, and left two
things unfixed. Some entries have \emph{no usable description}, so there is
nothing for a text tower to encode --- and a text index cannot rank a document
that does not exist, so those entries score exactly 0\%, structurally rather
than statistically. And some queries are \emph{ambiguous}, where a ranking is a
poor answer to a question the user has not finished asking.

Both are repaired by generation. A multimodal model can write the missing
description; and a generator with the retriever attached can answer the
ambiguous query in words, grounded in what the catalogue actually contains.

\section{Filling the holes}

60 of the 200 entries have no description. The captioner is architecturally the
encoder--decoder of Chapter 15 with an image encoder in place of the text one
--- a vision transformer, a transformer decoder cross-attending to it, trained
on image--caption pairs by next-token prediction, 224 million parameters.
Nothing new; a different pair of modalities in the same box. Nine lines, and 21
seconds for 60 images with beam search --- not instant, so run it offline, once
per new entry.

\begin{center}
\small
\begin{tabular}{p{0.42\textwidth}p{0.42\textwidth}}
\toprule
\textbf{Human description (deleted)} & \textbf{Generated description} \\
\midrule
This wire metal rack holds several pairs of shoes and sandals &
a dog laying on top of a pair of shoes \\
A green vase filled with red roses sitting on top of table &
a vase of flowers sitting on a window sie \\
A giraffe bending over while standing on green grass &
the gi is brown and white \\
\bottomrule
\end{tabular}
\end{center}

The first row is the interesting one: \emph{both sentences are true}. There is
a dog asleep on the shoe rack, and the two writers chose different subjects.
The third row is not a sentence at all.

\begin{center}
\small
\begin{tabular}{lrr}
\toprule
\textbf{Text-only route, all 200 queries} & \textbf{R@1} & \textbf{R@10} \\
\midrule
Every entry described by a person & 61.5\% & 96.0\% \\
60 entries with nothing & 46.0\% & 68.5\% \\
Those 60 auto-captioned & 61.5\% & 94.0\% \\
\bottomrule
\end{tabular}
\end{center}

At rank 1 the recovery is exact --- the auto-captioned index returns the same
61.5\% as the one every person wrote, which is a tie rather than a win. At rank
10 it is \emph{not} complete, and that gap is the real finding.

\begin{keybox}[title={Meanwhile, the route that never read a description}]
On the same 60 entries, the text index with descriptions deleted scores 0\%,
the auto-captioned index recovers most of it, and the joint-space image route
--- which never needed text at all --- scores 80.0\%.

\emph{Two routes that fail for different reasons are worth more than one better
route.} The image route is immune to missing text; the text route is immune to
a photograph that does not show the distinguishing feature.
\end{keybox}

\begin{failbox}[title={Failure condition --- three things not to claim}]
A generated caption is \emph{not evidence about the product}. It is a
restatement of the photograph, and it will confidently name things that are not
there.

The recovery is measured on 60 entries of 200 --- a small sample with a wide
interval. And we are scoring generated captions against human captions of the
same image, which is a friendly test: a customer's wording is not a COCO
caption.

The failure mode to watch for: an auto-caption that is \emph{wrong} makes an
entry findable under the wrong query, which is worse than unfindable. Nothing
we measured detects that.
\end{failbox}

So the honest use is narrow and worth stating: fill holes, mark the filled
entries as machine-written, and measure retrieval with and without them.

\begin{keybox}[title={Prompting is a hyperparameter}]
A generator's output depends on how you ask, and the difference is often larger
than the difference between models. Trying five prompts and keeping the one
with the best test score is choosing a hyperparameter on the test set ---
Lecture 2's error, in a costume that does not look like modelling at all.

Report the prompt with the result. A number without the prompt that produced it
is not reproducible, and prompts are rarely written down because they do not
feel like code.
\end{keybox}

\section{Retrieval-augmented generation}

The user asks in words; the retriever of Lecture 20 embeds the question and
returns the $k$ most similar entries; those go into the generator's context
with the question; and the generator answers \emph{from what it was given},
citing which entry it used.

Why this is not just a bigger model: \emph{the knowledge lives in the index,
not in the weights}. Add a product and it is answerable immediately; remove one
and it stops being cited. No retraining --- and the system can say where an
answer came from, which is usually the requirement that decides whether it can
be deployed at all.

\begin{center}
\small
\begin{tabular}{lp{0.30\textwidth}p{0.34\textwidth}}
\toprule
& \textbf{Good at} & \textbf{Bad at} \\
\midrule
Retriever & finding the right entries, fast, over millions &
            saying anything; it only ranks \\
Generator & fluent answers, synthesis across entries &
            knowing anything reliably; it will produce a plausible answer
            regardless \\
\bottomrule
\end{tabular}
\end{center}

\begin{failbox}[title={Failure condition --- visible from neither half}]
\emph{Retrieved but ignored} --- the right entry is in the context and the
generator answers from its own weights instead.

\emph{Not retrieved, answered anyway} --- the retriever misses and the
generator does not say so, producing a fluent answer with nothing behind it.

\emph{Retrieved and misread} --- the entry is present and correct, and the
answer misstates it.

Evaluating the retriever alone by recall@$k$ and the generator alone by fluency
can both look excellent while the system does all three. \emph{The join needs
its own measurement.}
\end{failbox}

Two upstream decisions are usually made without measurement. \emph{Chunking}:
too small and a chunk lacks the context that makes it answerable; too large and
the embedding averages several topics and matches none well; split badly and a
sentence cut in half retrieves for neither of its halves. It is a
hyperparameter --- vary it, hold everything else fixed, and measure recall@$k$
on the same judgements. And \emph{how many to retrieve}: small $k$ risks the
right entry not being there at all, large $k$ buries it among distractors ---
which is where re-ranking earns its cost, retrieving generously with the cheap
bi-encoder and re-ranking with the cross-encoder, high recall from the first
stage and high precision from the second.

\subsection{Decide how you will check it first}

``The answer is good'' is not measurable in a lecture. So give every entry a
stock number and require the assistant to cite them: a cited SKU either exists
in the catalogue or it does not. Count the citations, count the ones that
exist, report the fraction --- no judgement call anywhere in the loop.

\emph{This measures grounding, not helpfulness.} Say which one you measured;
they are not the same thing and only one of them is cheap.

\begin{center}
\small
\begin{tabular}{lrr}
\toprule
\textbf{Over 12 ambiguous queries} & \textbf{Closed book} &
\textbf{Retrieval-augmented} \\
\midrule
Queries that cited at least one SKU & 10 & 12 \\
Citations emitted & 32 & 24 \\
Citations that exist & 0 & 24 \\
Grounding rate & 0\% & 100\% \\
\bottomrule
\end{tabular}
\end{center}

Closed book, the model produces a fluent recommendation and cites stock numbers
\emph{in exactly the right format}. Across twelve queries it emitted 32
citations, of which none corresponds to an entry that exists.

\begin{keybox}[title={The failure is not that it refuses}]
It is that it does \emph{not} refuse, and nothing in the output distinguishes
an invented SKU from a real one.
\end{keybox}

The generator is a 494-million-parameter instruction model, chosen so the whole
thing runs on a laptop; a larger one would follow the citation instruction more
reliably and would not change the shape of the result.

\begin{failbox}[title={Failure condition --- what this does not fix}]
Grounding is not correctness --- a cited entry can exist and still be a bad
recommendation. We measured the cheap half.

The retriever is now the ceiling: if the right entry is not in the top five, no
amount of generation recovers it, which is why the R@5 measured last lecture is
the number that matters here.

And fluency is unchanged. A wrong answer built from real SKUs reads better than
a wrong answer built from invented ones.
\end{failbox}

Four ways to measure the join: \emph{groundedness}, is every claim supported by
a retrieved entry; \emph{attribution}, does the cited entry actually say it ---
a citation that does not support its claim is worse than none;
\emph{abstention}, when nothing relevant is retrieved does it decline, measured
deliberately with queries whose answer is genuinely absent; and only then, end
to end, is the answer right. Plus the closed-book control: run the same
questions with retrieval off, and whatever the system still answers correctly
it knew already. \emph{The difference is what retrieval actually bought} ---
the same shape as every ablation in this course.

\begin{center}
\small
\begin{tabular}{lrl}
\toprule
\textbf{Route} & \textbf{Before} & \textbf{After} \\
\midrule
Text query to image, joint space & 73.0\% & 73.0\% (untouched) \\
Text index, on the entries with no description & 0\% & recovered \\
Ambiguous queries, grounding rate & 0\% & 100\% \\
\bottomrule
\end{tabular}
\end{center}

Two repairs, two measured deltas, and \emph{one number deliberately unchanged
as a control}. That last row is what makes the other two believable.

\section{The course, as one argument}

\begin{keybox}[title={What the first lecture promised}]
``How to build a machine learning system, and how to know whether it works.
Fitting a model is a few lines. Establishing that the number it reports means
what it appears to mean is the whole of the rest of the course.''

That was a claim, not a plan. The claim, stated precisely: \emph{every method
you met was introduced at the moment a number you produced was wrong without
it.} Not ``was useful''. Not ``is standard''. Was wrong, and measurably so, and
you were the one holding the wrong number.
\end{keybox}

\subsection{The eighteen derivations, and where each returned}

\begin{center}
\small
\begin{tabular}{lp{0.30\textwidth}p{0.34\textwidth}}
\toprule
\textbf{Lecture} & \textbf{The derivation} & \textbf{Where it came back} \\
\midrule
2 & least squares, the normal equation & ridge in 5; collinearity everywhere \\
3 & imbalance; precision is not monotone & the maximum inside AP, Lecture 14 \\
4 & gradient descent & every network from Lecture 9 on \\
5 & bias--variance & every learning curve since \\
6, 7 & impurity; variance of a correlated average & multi-head attention,
       Lecture 18 \\
8 & PCA via the SVD; Johnson--Lindenstrauss & factorisation in 21; the sphere
    in 23 \\
9 & what a layer computes & everything after it \\
10 & reverse-mode autodiff & why training is affordable at all \\
11 & variance propagation & BPTT in 16; residuals in 18 \\
12 & weight sharing, equivariance, memory & what a ViT gives up, Lecture 23 \\
14 & IoU; mAP & ranking metrics in 19 \\
15 & stationarity; autocorrelation & causal masking in 18 \\
17 & softmax, cross-entropy, logits & attention's softmax in 18; the
     contrastive loss in 23 \\
18 & scaled dot-product attention & Part V's encoders, and 23 \\
19 & MRR, AP, NDCG & a ranking needs its own metrics \\
21 & matrix factorisation and the SVD & Lecture 8, with missing entries \\
23 & the contrastive objective and its temperature & the same two towers,
     different modalities \\
\bottomrule
\end{tabular}
\end{center}

Lectures 20, 22 and 23 are three uses of \emph{one object} --- two encoders and
an inner product. That is the strongest single thread in the course, and it is
why Part V sits after transformers rather than before.

\subsection{The failures, collected}

\begin{center}
\small
\begin{tabular}{llp{0.36\textwidth}}
\toprule
\textbf{Failure} & \textbf{First met} & \textbf{Measured cost} \\
\midrule
Fitting before splitting & Lecture 2 & nothing measurable --- and unknowable in
                                       advance \\
Evaluating on training data & Lecture 2 & an RMSE of zero \\
Accuracy under imbalance & Lecture 3 & a useless model above 90\% \\
Choosing on the test set & Lectures 2, 5 & an optimistic report, every time \\
Missing \code{zero\_grad()} & Lecture 10 & 76.7 points \\
Default initialisation, deep & Lecture 11 & a network that does not train at
                                            all \\
Random split on a series & Lecture 15 & optimistic, and undeployable \\
Augmenting the validation set & Lecture 13 & 1.57 points of noise on the
                                             criterion \\
Embeddings compared without normalising & Lecture 23 & a different ranking \\
\bottomrule
\end{tabular}
\end{center}

\emph{Not one of them raises an exception.} That is the whole reason the course
is organised around measurement.

Two of them are worth restating with their arithmetic. In detection, the same
detections at the same IoU threshold give 0.659 when averaged over classes and
0.817 when averaged over images --- \emph{a mean of a mean is not a mean}, and
the headline moves by 0.157 depending on which you pick. Against the stricter
0.5-to-0.95 average of 0.439 the apparent gap looks like 0.378, and that
comparison is not a comparison at all, because the two rows are averaged over
different thresholds.

And in forecasting: the same model on the same rows scores 44{,}925 under a
random split and 56{,}446 under a time-ordered one --- 25.6\% optimism, with
nothing fitted on the test set and no preprocessing crossing the split. The
leak was \emph{the ordering}, and with positive autocorrelation a point's own
near-future was in its training set.

\begin{keybox}[title={What did not change, in twenty-four lectures}]
\begin{enumerate}[leftmargin=1.6em,itemsep=1pt]
  \item Split before anything is fitted.
  \item Compute the trivial baseline first --- a metric with nothing beside it
        is decoration.
  \item Two numbers measured on different rows are not comparable, however
        similar the column headings look.
  \item A difference smaller than the spread is not a difference. Report the
        spread.
  \item Change one thing, or the difference is not attributable to anything.
\end{enumerate}
Five rules, no mathematics, and they would have caught every failure in the
table above.
\end{keybox}

\subsection{What this course did not cover}

``I took a machine learning course'' is read by other people as a claim about
what you can do, so you should be able to say precisely what it does and does
not mean. The scope was Chapters 1--16 of one book, thoroughly, on thirteen
datasets. That is a large and useful slice of applied machine learning. It is
not the field.

Missing, and each is a course: reinforcement learning, a different problem
shape entirely; generative image models; causality --- every method here
predicts, and none tells you what happens if you \emph{intervene}; fairness and
model governance; and serving and drift.

\begin{keybox}[title={The two you are most likely to need first}]
\emph{Drift.} Every number in this course was measured on data drawn like the
training data. In production that stops being true, gradually and without
warning, and the score you reported at launch is the last honest one anybody
has. Monitor the inputs, not just the outputs --- the distribution moves before
the accuracy does.

\emph{Who it is worse for.} Lecture 2 sliced the error by district and found
the poorest tenth predicted worst. That slice was one line of code, and it is
the difference between a system you can defend and one you cannot. Always ask
for whom the average is wrong.
\end{keybox}

And the part that moves fastest: pretraining at scale --- we used four
pretrained models in this last application and trained none of them;
instruction tuning, which is why the generator behaves as it does; and
evaluating generated text, where we measured grounding \emph{because grounding
is checkable} and measuring whether an answer is good is unsolved.

\begin{keybox}[title={The transferable part}]
Every one of those still fails in the ways this course taught you to look for:
the wrong metric, a missing baseline, a test set that leaked, and a number
quoted from one run.
\end{keybox}

Without a ground truth you can still measure \emph{consistency} --- ask the
same thing twice in two wordings, and disagreement bounds reliability from
above; \emph{groundedness} against the retrieved text, mechanically;
\emph{abstention} on deliberately unanswerable queries, the cheapest useful
test of a RAG system; and human judgement on a sample, expensive and still the
anchor everything else is validated against.

\section{What you can do now}

Take a problem, frame it, and say what a correct answer would be before fitting
anything. Choose a metric that matches the brief, and measure the trivial
baseline first. Build the simplest thing that runs, inside a pipeline,
evaluated on a protocol fixed in advance. Read the failure --- slice the error,
look at the worst cases, name the conditions under which the method breaks.
Improve one thing at a time and report the ablation rather than the best row.
And write down what the number was measured on, every time.

Six steps. Nothing in them is specific to a library, and none will be out of
date in ten years.

\begin{keybox}[title={A last word on the assistant}]
You will build most of this with an assistant, and Lecture 1 set out how:
specify, generate, read, test, verify --- and steps three to five are yours.

What has changed by now is that you know what to look for. ``What was this
measured on?'', ``what is the baseline?'', ``was anything fitted before the
split?'' are the questions an assistant will not ask itself, and they are
precisely the ones that decide whether the number it hands you means anything.

\emph{The code was never the scarce part. Knowing whether to believe the output
is, and that is what these twenty-four lectures were for.}
\end{keybox}

How to keep learning: reproduce something --- take a paper's headline number
and try to get it, and you will learn more from the gap than from the paper.
Keep the habit of the baseline. Read the evaluation section first: if it does
not say what was held out and how, the rest is decoration. And write the check
before the model.

\emph{The libraries in these notebooks will be superseded. The derivations will
not, and neither will the habit of asking what a number was measured on.}

\section{The notebook}

\begin{trybox}[title={What to change}]
Ask a question whose answer is not in the catalogue. A system that answers it
anyway has the second failure of the join --- and you have just built the test
that finds it.
\end{trybox}

\section{Exercises}

\begin{enumerate}[leftmargin=1.6em,itemsep=4pt]
  \item An entry with no description scores 0\% on a text retrieval route.
        State why this is a structural rather than a ranking failure, and give
        two independent repairs. \hfill \scope{[4]}
  \item Name the three failures that belong to the join of a retriever and a
        generator, and explain why evaluating each half separately cannot find
        them. \hfill \scope{[5]}
  \item A retrieval-augmented system is judged by a grounding rate. State
        precisely what that measures, what it does not, and the control you
        would run beside it. \hfill \scope{[5]}
  \item The same detections give mAP 0.659 averaged over classes and 0.817
        averaged over images. Explain the discrepancy, and say which number you
        would report and why. \hfill \scope{[4]}
  \item Give the five procedural rules this course did not change, and for each
        name one failure it would have caught. \hfill \scope{[5]}
\end{enumerate}

\section*{Summary}
\addcontentsline{toc}{section}{Summary}

Generation repaired the two holes Lecture 23 left. Captioning took entries that
were structurally invisible to a text index --- not badly ranked, absent --- and
made them findable, recovering the rank-1 figure to within one query in 200 of
what human writers achieved and \emph{not} recovering rank 10, which is the
finding rather than the headline. Meanwhile the image route, which never needed
a description at all, scored 80.0\% on those same entries: two routes that fail
for different reasons are worth more than one better route.

Retrieval-augmented generation took twelve ambiguous queries from a 0\%
grounding rate to 100\%. Closed book, the model emitted 32 citations in exactly
the right format, of which none existed --- and the failure is not that it
refuses but that it does not, with nothing in the output distinguishing an
invented identifier from a real one. Grounding is not correctness; we measured
the cheap half, and said so.

And a control row deliberately left unchanged is what makes the two measured
deltas believable.

The course was one argument: every method arrived at the moment a number was
wrong without it. Eighteen derivations, thirteen datasets, and a catalogue of
failures none of which raises an exception --- a scaler fitted early that cost
nothing measurable, a missing \code{zero\_grad()} that cost 76.7 points, a
random split on a series that cost 25.6\% and deployability. The five
procedural rules that would have caught all of them contain no mathematics at
all.

\end{document}
